\section{Method}
\label{sec:method}

\subsection{Tight GBDN: nonsubsampled analysis}
\label{sec:method-tight}

At level $\ell$, fix a root multiset $\mathcal R_\ell$ independently of the
input being analyzed and define
\begin{equation}
  P_\ell^{+}=\frac{I+T_\ell}{2},
  \qquad
  P_\ell^{-}=\frac{I-T_\ell}{2},
  \qquad T_\ell=\Tcal_{\mathcal R_\ell}.
  \label{eq:tight-channels}
\end{equation}
Starting from a complex feature signal $h_0$, Tight GBDN emits and carries
\begin{equation}
  r_\ell=P_\ell^-h_\ell,
  \qquad
  h_{\ell+1}=P_\ell^+h_\ell,
  \qquad \ell=0,\ldots,D-1.
  \label{eq:analysis-recursion}
\end{equation}
The complete coefficient representation, in the canonical order used by the
implementation and artifacts, is
\begin{equation}
  \Acal_Dh_0=(r_0,\ldots,r_{D-1},h_D).
  \label{eq:analysis-map}
\end{equation}
This is a redundant, nonsubsampled analysis bank: every channel lives on all
$n$ vertices.  We use ``pointwise paraunitary'' only for the multilevel
scalar spectral analysis vector whose Hermitian squared norm is one
(\cref{thm:pointwise-partition}); the term does not imply downsampling,
polyphase structure, or critical sampling.

A predictor concatenates real and imaginary parts of the coefficients.  Its
learned lift and readout are outside the analysis guarantee.

\paragraph{Two reconstruction identities.}
For coefficients produced from a common state by the same channel pair,
additive reconstruction is the algebraic identity
\begin{equation}
  h_\ell=r_\ell+h_{\ell+1},
  \label{eq:additive-reconstruction}
\end{equation}
which holds because $P_\ell^-+P_\ell^+=I$.  For the finite realization define
$\widetilde P_{\ell,K}^{\pm}=(I\pm\widetilde T_{\ell,K})/2$ and the analyzed
coefficients, starting from $\widetilde h_0=h_0$,
$\widetilde r_\ell=\widetilde P_{\ell,K}^-\widetilde h_\ell$ and
$\widetilde h_{\ell+1}=\widetilde P_{\ell,K}^+\widetilde h_\ell$; then
$\widetilde r_\ell+\widetilde h_{\ell+1}=\widetilde h_\ell$ as well.
When $K$ is fixed, we suppress it in $\widetilde T_\ell$ and
$\widetilde P_\ell^\pm$.

Adjoint synthesis is distinct.  For an exact analyzed tuple it is the
backward recursion
\begin{equation}
  h_D^{\mathrm{syn}}=h_D,
  \qquad
  h_\ell^{\mathrm{syn}}
  =(P_\ell^-)^*r_\ell+(P_\ell^+)^*h_{\ell+1}^{\mathrm{syn}},
  \quad \ell=D-1,\ldots,0.
  \label{eq:synthesis-recursion}
\end{equation}
It yields $h_0^{\mathrm{syn}}=\Acal_D^*\Acal_Dh_0=h_0$ because the exact
analysis is Parseval.  The finite counterpart is
\begin{equation}
  \widetilde h_D^{\mathrm{syn}}=\widetilde h_D,
  \qquad
  \widetilde h_\ell^{\mathrm{syn}}
  =(\widetilde P_{\ell,K}^-)^*\widetilde r_\ell
   +(\widetilde P_{\ell,K}^+)^*
     \widetilde h_{\ell+1}^{\mathrm{syn}},
  \quad \ell=D-1,\ldots,0.
  \label{eq:approx-synthesis-recursion}
\end{equation}
It returns
$\widetilde h_0^{\mathrm{syn}}=
\widetilde\Acal_D^*\widetilde\Acal_D\widetilde h_0$ and need not be an
inverse.  Its error guarantee applies only to coefficients produced by the
same approximate analysis, not to an arbitrary coefficient tuple.  Neither
reconstruction identity implies that the preceding learned feature lift is
invertible or that the original real input is recoverable.

\subsection{Product-sum GBDN}
\label{sec:method-product-sum}

The expressive variant defines cumulative products
\begin{equation}
  Q_0=I,
  \qquad
  Q_\ell=T_\ell Q_{\ell-1},
\end{equation}
and a learned complex scalar-weighted sum
\begin{equation}
  H_{\mathrm{PS}}h=\sum_{\ell=0}^{D}c_\ell Q_\ell h,
  \qquad c_\ell\in\C.
  \label{eq:product-sum}
\end{equation}
Each exact factor is unitary, but the learned sum is not generally unitary or
tight.  This is a BDN-inspired secondary model, not an alternative statement
of the Tight GBDN guarantees.

\subsection{GBDN+: empirical relaxation}
\label{sec:method-relaxed}

GBDN+ applies parallel root-parameterized filters to the same lifted input,
adds unconstrained Chebyshev corrections, and learns a skip mixture.  The
corrections remove exact all-pass, tightness, and reconstruction guarantees.
We therefore evaluate canonical GBDN+, and label the historical implementation
Legacy GBDN+, only as empirical architectures; neither inherits the theorems
for Tight GBDN.

\subsection{Finite sparse realization and graph identity}
\label{sec:method-cheb}

All sparse variants evaluate \cref{eq:cheb-realization} by recurrence.  A
caller may pass a validated self-adjoint operator for a fixed graph; otherwise
the implementation constructs the normalized Laplacian from the supplied
graph identity.  It never reuses an operator merely because two graphs share
the same vertex count.  Under the canonical full recurrence, each degree-$K$
factor is evaluated without early truncation or cross-level recurrence reuse
and therefore performs exactly $K$ sparse operator applications.  One
application of the real Laplacian to a complex $n\times d$ feature matrix is
counted as one complex sparse matrix--feature product (two products under a
real-channel count).  A different batching or truncation schedule must report
its own count.  The finite realization is the object used by GPU experiments;
exact operators serve as a dense spectral oracle on small graphs.
